\section{Conclusion intermédiaire}

Cette première étape installe la base méthodologique du projet :
\begin{itemize}
    \item une implémentation autonome de l'étape 1 conservée dans \texttt{stages/01\_numba\_threads/} ;
    \item une stratégie de benchmark reproductible, séparant calcul et affichage ;
    \item un rapport LaTeX structuré dès le début pour intégrer progressivement les résultats validés.
\end{itemize}

Sur la configuration retenue, le benchmark passe d'un temps moyen de \SI{227.970}{ms} par pas à \SI{33.233}{ms} avec 16 threads, soit un speed-up de \(6.861\). Cette étape valide donc la pertinence de l'approche Numba sur la partie calcul, tout en montrant déjà les limites d'efficacité attendues lorsque le nombre de threads augmente.
